% 结构版本 v0.2：用户已确认审阅稿（D015）；3.1 定稿 V015（2026-09-26）；3.2 骨架待第四至六章符号收全后随 TBL-003 起草。
\section{模型假设与符号说明}\label{sec:assumptions-notation}

\subsection{模型假设}\label{sec:assumptions}
% 状态：V015 定稿（2026-09-26 用户对话审定；四条候选预览稿经一轮 GPT 结构评审逐条仲裁后收口为两条制冻结，AI-16/A052）。素材：题面 B.4/C.1/B.6、附录 D.3--D.5、1.2.1/1.2.2/1.3.1/1.4/4.2.4 定稿回指锚点、三篇基准假设节写法对照（2025A 五条均为主动简化型——其确有可简化对象；2021A/2023A 短条目改变适用域型）。仲裁要点：①四条候选（输入合法性/核间资源共享范围/DDR 容量/评估确定性）被裁决信息类型错位——题面事实与官方评估器性质不列假设：输入合法性删除（1.2.2 题设保证复述＋1.3.1 已有载体、检验句属伪检验）、评估确定性移出（C.1"保证结果可复现"为官方性质）；②候选稿"跨核耦合全部经共享带宽资源表达"技术错误撤出（4.1 激活等待 100/1000、5.1 同步延迟 500、4.2.4 L_relay 逐段等待均为带宽外跨核耦合反例）；③DDR 条定性为官方评估抽象下的建模约定、不作物理容量断言（条目称"约定"）；④末句范围收口"执行性能与优化效果的结论"（图结构定理/复杂度等结论不来自评估器）；⑤机制枚举按 B.4 分场景（场景 A 的 Task 等待／场景 B 与问题三的跨核 COPY 同步延迟）；⑥"跨核通信"改"跨 Task 数据搬运"（对齐 1.4 新增搬运三来源口径，覆盖场景 A 同核跨 Task 情形）。Needs-integration：4.2.4 待补确定性半句（用户已批口径：对固定测试用例、核心数量与问题配置，同一合法方案经官方评估程序得到唯一的评价结果，可将其作为方案到评价指标的确定性映射——随 4.2.4 修订窗，不另占 3.1）。
本题的计算图、硬件配置与执行评价规则均由题面及官方评估程序给定，本文以官方评估语义作为建模与方案评价的统一依据，不对题面规定的依赖关系、缓存容量、共享带宽、同步机制及合法性约束作额外简化（见~\ref{sec:q1-evalbound} 节）。对于题面未显式规定、但为明确模型边界所必需的部分，作如下建模约定与假设。

\begin{itemize}
\item \textbf{DDR 存储空间不构成容量约束}。题面给出的存储容量参数仅涉及核内 L1、UB 与问题三的共享 L2，官方评估程序亦未设置 DDR 存储容量与地址分配约束；据此，本文按官方评估抽象处理，不对写入 DDR 的数据设置空间限制，不计核外存储碎片与地址管理开销，由缓存换入换出及跨 Task 数据搬运产生的 DDR 代价仅按官方规则体现为总额外数据搬运量、共享带宽占用与相应执行时间。该约定对三个问题一致，其依据为官方评估程序的配置项与执行逻辑核验。
\item \textbf{忽略官方评估模型之外的微架构效应}。除题面明确规定的执行管道、L1/UB 容量、DDR 共享带宽、场景 A 的 Task 等待、场景 B 与问题三的跨核 COPY 同步延迟，以及问题三的共享 L2 等机制外，本文不引入总线仲裁、寄存器冲突、信号传播与时钟抖动等未在官方评估程序中建模的硬件效应；本文报告的执行时间与优化效果均针对题面规定的 NPU 抽象模型，不直接等同于具体物理芯片的实测性能。该假设限定模型的适用范围，本文有关执行性能与优化效果的结论均以官方评估程序为评价依据，不外推至该抽象之外。
\end{itemize}

\subsection{符号说明}\label{sec:notation}
\pending{计算图、节点与边。 子图、核心及执行顺序。 时间、通信、存储与缓存状态。 目标函数与评价指标。 只收录正文实际使用的符号，统一定义域、单位及索引。}
